\chapterimage{images/15-posledice.jpg}

\chapter{Posledice}

\editorialnote{Sudbine važnih likova i most ka svetu Endimiona.}

Šest meseci posle smrti Mreže, Hiperion je izgledao mirnije nego što bi iko mogao da očekuje.

Rat u njegovom sistemu završio se ubrzo nakon pada dalekobacačke mreže.

Proterani su zadržali snažno prisustvo, ali nisu uništili stanovništvo niti pretvorili planetu u okupiranu pustinju.

U Kitsu je uspostavljena zajednička uprava sa preostalim predstavnicima Hiperiona.

Konzul i Teo Lejn pomogli su da se postigne sporazum.

Po ulicama prestonice sada su se mogli sresti ljudi Hegemonije i visoke, prilagođene prilike Proteranih.

Nekada bi prizor izazvao paniku.

Sada je postajao svakodnevan.

\bigskip

Stari poredak nije preživeo.

Ali ljudi jesu.

To nije značilo da je sve bilo dobro.

Saobraćaj između svetova gotovo je nestao.

Brodovi sa Hokingovim pogonom ponovo su bili jedina prava veza između zvezda.

Putovanje je ponovo značilo mesece ili godine za putnika i mnogo veći vremenski dug za one koje ostavlja iza sebe.

Trgovina je opala.

Svetovi su morali da uče kako da žive od sopstvenih resursa.

Mnoge stvari koje su nekada bile obične postale su luksuz.

Neke su potpuno nestale.

Ali najgora predviđanja o trenutnom kraju civilizacije nisu se ostvarila.

Ljudi su se prilagođavali.

Sporo.

Bolno.

Ali jesu.

\bigskip

Hiperion je u tom novom svetu ponovo postao ono što je bio pre nego što ga je Hegemonija pokušala da uključi u Mrežu.

Udaljen.

Neobičan.

Teško dostupan.

Samo što sada gotovo svaki ljudski svet deli nešto od njegove nekadašnje usamljenosti.

\bigskip

Bron Lamija ostala je u Kitsu.

Sedmi mesec trudnoće već joj je menjao hod, ali nije promenio karakter.

Stanovala je kod Cicerona.

Sten Leveski je obnavljao oštećenu gostionicu i svakodnevno joj govorio da trudnice ne treba da nose daske, pomeraju nameštaj ni pomažu radnicima.

Bron ga uglavnom nije slušala.

Džoni je bio mrtav.

Njegova ličnost iz Šrenove petlje bila je uništena.

Ali njegovo dete je živelo.

I Bron više nije mogla da tvrdi da je to samo ostatak jednog izgubljenog slučaja.

Bila je to njena kćer.

\bigskip

Jedne večeri ponovo je ugledala Džona Kitsa.

Ne Džonija.

Drugu rekonstrukciju.

Severna.

Ili ono što je od njega ostalo.

Više nije imao telo.

Njegova ličnost mogla je kratko da se pojavi tamo gde su još postojala dovoljno snažna polja metasfere.

Bron je posegnula da ga dodirne.

Ruka joj je prošla kroz njega.

``Ti si drugi Kits.''

Nasmešio se.

``Može se tako reći.''

Bron je spustila pogled na stomak.

``Ti si spasao Rejčel.''

``Na kratko.''

``Zašto?''

Kits je pogledao njeno nerođeno dete.

``Zato što je trebalo da je neko sačeka.''

Bron ga je proučavala.

``Jesi li ti Empatija? Onaj deo ljudske Ultimativne Inteligencije koji je pobegao kroz vreme?''

``Ne.''

Njegov odgovor bio je neposredan.

``Ja sam samo Onaj Koji Dolazi Ranije.''

Bron se namrštila.

``Šta to znači?''

``Da pripremam put.''

``Za koga?''

Kits je ponovo pogledao njen stomak.

``Za Onu Koja Podučava.''

\bigskip

Bron je sela.

``Moju kćer?''

``Tvoju i Džonijevu.''

``Šta će ona podučavati?''

``Ne znam.''

Bron ga je pogledala sumnjičavo.

``Odlično. Proročanstvo bez informacija.''

Kits se nasmešio.

``Znam samo da će njene ideje menjati svetove veoma dugo.''

Bron je položila ruku na stomak.

Prvi put nije pokušala da pronađe racionalno objašnjenje.

Nije ga bilo.

Još ne.

\bigskip

Kits joj je rekao još nešto.

Kada se na Hiperionu suočila sa Šrajkom i učinila ono što ni sama nije razumela, nije joj pomogla Moneta.

Nije joj pomogao ni on.

Ono što je učinila poticalo je od nje same.

Ili od nečega što je već bilo povezano s detetom koje nosi.

Bron je odmah odbacila tu mogućnost.

Kits se samo osmehnuo.

``Majka One Koja Podučava možda ima nekoliko privilegija.''

``Ne počinji.''

\bigskip

Njegovo prisustvo je slabilo.

Srž više nije postojala na način na koji je nekada postojala.

Datasfere pojedinačnih svetova bile su premale da dugo nose njegovu ličnost.

Postojala je metasfera.

Ali Kits još nije bio spreman da ode tamo.

Bron nije pitala šta ga čeka.

Po izrazu njegovog lica zaključila je da ni on ne zna.

``Hoćeš li se vratiti?''

``Možda.''

``Mrzim takve odgovore.''

``Znam.''

Kada je ponovo podigla pogled, više ga nije bilo.

\bigskip

Fedman Kasad nije se vratio živ.

U dalekoj budućnosti stao je pred vojsku ljudi koja je čekala poslednju bitku protiv legija Šrajkova.

Moneta ga tada još nije poznavala kao čoveka kojeg će voleti.

Za nju je njihov odnos tek trebalo da počne.

Za Kasada je trajao gotovo čitav život.

\bigskip

Kada je došao trenutak za borbu, Kasad više nije ratovao za Hegemoniju.

Ni za SILU.

Ni za slavu.

Borba se vodila za pravo čovečanstva da odlučuje o sopstvenoj prošlosti i budućnosti.

Kasad je poveo napad.

I poginuo.

\bigskip

Posle bitke pronađen je u smrtnom zagrljaju sa Šrajkom protiv kojeg se borio.

Njegovo telo je očišćeno, obučeno u crnu uniformu SILE i položeno u Kristalni Monolit.

Oružje mu je ostavljeno kraj nogu.

Moneta se oprostila od njega.

Zatim je ušla u Sfingu i nastavila sopstveno putovanje kroz vreme.

\bigskip

Kasad je tako postao deo upravo one legende kojoj je čitavog života nesvesno išao u susret.

Njegova grobnica stajala je na Hiperionu vekovima pre nego što je on u njoj sahranjen.

\bigskip

Het Mastin nije dobio takvu drugu priliku.

Njegova misija postala je jasna tek nakon njegove smrti.

Templari su znali da će Igdrasil biti uništen.

Mastin se dobrovoljno prihvatio uloge koju je predviđalo njihovo proročanstvo.

Trebalo je da na Hiperionu primi Drvo Ispaštanja.

Šrajkovo Drvo bola.

Uz pomoć erga trebalo je da postane njegov Glas i da ga vodi kroz prostor i vreme.

Kada je zaista ugledao Drvo, nije mogao.

Možda se uplašio.

Možda je odbio.

Možda je prvi put shvatio da proročanstvo njegovog reda nije sveto onako kako je verovao.

Kakva god bila istina, nije izvršio zadatak.

Umro je pre nego što je mogao potpuno da objasni svoj izbor.

Njegova priča ostala je jedina hodočasnička priča koja nikada nije ispričana njegovim glasom.

\bigskip

Lenar Hojt takođe je ostao mrtav.

Ali kruciforma je njegovom smrću vratila Pola Direa.

Dire je završio daleko od Hiperiona.

Na Pacemu.

U svetu Katoličke crkve koja je i sama jedva preživela pad Mreže.

A onda mu je saopšteno nešto što je smatrao toliko besmislenim da je u početku odbio da poveruje.

Kardinali su ga izabrali za Papu.

\bigskip

``Mene?''

Dire je gledao poruku.

``Mora da postoji greška.''

Nije je bilo.

Čovek koji je bio ekskomuniciran.

Čovek koji je falsifikovao arheološke podatke.

Čovek koji je umro na tesla-drvetu na Hiperionu.

Čovek koji se vratio kroz tuđe telo.

Sada je trebalo da vodi jednu od najstarijih ljudskih religija.

Dire nije znao da li je to znak Božje volje ili još jedna kosmička šala.

U svakom slučaju, više nije mogao da se vrati životu koji je imao.

\bigskip

Martin Silenus se vratio sa Drveta bola.

I još jednom preživeo nešto što je trebalo da ga ubije.

Njegov \emph{Spev o Hiperionu} bio je završen.

Vekovima je verovao da mu je Šrajk potreban kao muza.

Da bez bola nema velike umetnosti.

Na Drvetu je dobio dovoljno bola za nekoliko života.

I otkrio da odgovor nije tako jednostavan.

Kada ga je Sol kasnije pitao za poemu, Martin je rekao:

``Dovršio sam je.''

Zatim je dodao:

``A otkrio sam i da pesnici nisu Bog.''

Sol je čekao nastavak.

``Ali ako Bog postoji, onda je verovatno pesnik.''

Martin je zastao.

``I prilično loš.''

Sol se nasmejao.

Rejčel je podrignula.

Martin je zaključio da je to možda najpoštenija kritika njegove teologije.

\bigskip

Najvažniji povratak bio je Rejčelin.

Sol je čekao pred Sfingom.

Nije se pomerio kada su stigli brodovi.

Nije se pomerio kada su se ostali vratili u dolinu.

Rat, Hegemonija i TehnoCentar više ga nisu zanimali.

Čekao je svoje dete.

\bigskip

Pred jutro se u svetlosti Sfinge pojavila žena.

U rukama je nosila bebu.

Sol ju je prepoznao istog trenutka.

Bila je to Rejčel.

Odrasla Rejčel.

A dete koje je nosila bila je takođe Rejčel.

Novorođenče koje je Sol predao Šrajku.

\bigskip

Odrasla kćer predala mu je svoju mlađu verziju.

``Da li je dobro?''

``Jeste.''

Rekla mu je da će beba od sada normalno stariti.

Merlinova bolest je završena.

Sol je privio novorođenče uz sebe.

A onda je pogledao odraslu ženu pred sobom.

I još jednom pokušao da razume kako obe mogu istovremeno postojati.

\bigskip

Bron je prva izgovorila ono što je ostalima promaklo.

``Moneta.''

Odrasla Rejčel ju je pogledala.

``Ti si Kasadova Moneta.''

Rejčel je potvrdila.

\bigskip

U budućnosti je odrasla u vremenu poslednjeg rata između ljudske i mašinske inteligencije.

Njena nekadašnja Merlinova bolest učinila ju je pogodnom za posebno putovanje kroz vreme.

Sa Vremenskim grobnicama krenula je unazad.

Postala je čuvar Šrajka.

I žena koju je Fedman Kasad kroz svoj život poznavao kao Monetu.

\bigskip

Za Kasada je njihova ljubav već bila prošlost.

Za Rejčel tek budućnost.

Videla ga je kako gine i ispratila njegovo telo u Kristalni Monolit, ali u svom ličnom toku vremena tek je trebalo da ga upozna.

Tako se njihova vremenska petlja zatvorila.

\bigskip

Odrasla Rejčel nije mogla dugo da ostane.

Morala je nazad u svoje vreme.

Pre nego što je otišla, rekla je Solu da i on može kroz Sfingu.

Sa bebom.

To bi značilo da će svoju kćer morati da odgaji još jednom.

Treći put.

``Nijedan roditelj to ne bi trebalo da mora'', rekla je.

Sol se nasmešio.

``Nijedan roditelj to ne bi odbio.''

\bigskip

Odrasla Rejčel vratila se u svetlost.

Sol je ostao sa bebom u naručju.

Ovoga puta dete je raslo u pravom smeru.

\bigskip

Sfinga se posle otvaranja promenila.

Više nije bila samo Grobnica.

Postala je prolaz kroz vreme.

Ne dalekobacač.

Nešto drugo.

Konzul nije znao kako da ga nazove.

Martin jeste.

``Vremebacač.''

Bron ga je udarila po ramenu.

``Nemoj da izmišljaš reči.''

``Kasno je za to.''

\bigskip

Proterani su mesecima proučavali temporalna polja.

Ljudi su pokušavali da uđu.

Većina nije mogla.

Sfinga je, izgleda, birala kada će nekoga propustiti.

Sol je odlučio da ode.

Sa Rejčel.

Ne nazad u stari život.

Taj život više nije postojao.

Sarai je bila mrtva.

Njegov univerzitet pripadao je drugom svetu i drugom vremenu.

Mreža je nestala.

Ali kroz Sfingu je postojala mogućnost da ode tamo gde je Rejčelin put zapravo vodio.

\bigskip

Ostali su ga ispratili.

Martin mu je pružio ruku.

``Priredi im pakao tamo, Vejntraube.''

``Kome?''

``Kome god da nađeš.''

Sol se nasmejao.

Martin je pogledao Rejčel.

``A ti pazi na starog.''

\bigskip

Bron je držala bebu dok se opraštala sa Solom.

``Je li sve u redu sa vašim detetom?'' pitao je.

``Jeste.''

``Dečak ili devojčica?''

``Devojčica.''

Sol se nasmešio.

Poljubio ju je u obraz.

Bron je skrenula glavu da drugi ne vide suze.

``Devojčice su strašno naporne'', rekao je Sol. ``Zamenite je za dečaka prvom prilikom.''

``Važi.''

\bigskip

Onda je uzeo Rejčel.

Ranac mu je bio na leđima.

Njegova kćer u naručju.

Pred njima je svetleo ulaz u Sfingu.

``Bilo bi prilično smešno da ovo sada ne proradi.''

``Proradiće'', rekao je Konzul.

``Kako znaš?''

``Ne znam.''

Martin je klimnuo.

``Dobar naučni princip.''

\bigskip

Sol se okrenuo prema njima.

``Zbogom.''

Napravio je korak.

Temporalna polja koja su ga nekada odbijala sada su ga propustila.

``Vredelo je'', rekao je.

Zatim su on i Rejčel nestali u svetlosti.

\bigskip

Nekoliko trenutaka niko nije govorio.

Hodočašće koje je počelo sa sedmoro ljudi završilo se tako da gotovo niko nije ostao tamo gde je počeo.

Hojt je bio mrtav.

Dire je ponovo živeo i išao prema životu koji će ga odvesti do papstva.

Mastin je bio mrtav.

Kasad je poginuo u dalekoj budućnosti i počivao u Kristalnom Monolitu.

Sol je sa novorođenom Rejčel otišao kroz Sfingu prema budućnosti u kojoj će je ponovo odgajati.

Bron je nosila Džonijevu kćer.

Silenus je preživeo Drvo bola i dovršio svoju poemu.

Konzul više nije imao Hegemoniju kojoj bi služio niti onu koju bi mogao da izda.

\bigskip

Ni Vremenske grobnice više nisu bile ono što su bile pre hodočašća.

Njihova polja su se otvorila.

Njihova svrha bila je samo delimično jasnija.

Sfinga je postala jednosmerni prolaz prema budućnosti.

Kroz nju su otišli Sol i Rejčel.

Kristalni Monolit postao je Kasadova grobnica.

Na njegovim oznakama govorilo se o velikom ratniku koji dolazi iz prošlosti da bi se borio protiv Boga Bola.

\bigskip

Druge Grobnice ostale su zagonetnije.

Neke su izgleda vodile prema udaljenim svetovima ili drugim vremenima.

Obelisk je ostao zatvoren.

Šrajkova Palata ostala je povezana sa mestima i vremenima koja ljudi još nisu potpuno razumeli.

Proterani i preostali stručnjaci SILE nastavili su da proučavaju dolinu.

Ali više niko nije mogao ozbiljno da tvrdi da su to samo stare ruševine.

\bigskip

Ni Šrajk nije bio konačno objašnjen.

Bron je uništila jednog.

Kasad se borio protiv drugog ili možda istog bića u drugom vremenu.

U budućnosti su postojale čitave legije Šrajkova.

Zato niko nije znao da li je opasnost nestala ili se samo povukla iz sadašnjosti.

Konzul je pogledao prema Grobnicama.

``Pretpostavljam da ćemo uvek imati Šrajka.''

Martin je slegnuo ramenima.

``Ili bar glasine o njemu.''

\bigskip

Šest meseci kasnije, Bron, Martin, Teo i Konzul sedeli su na balkonu Konzulovog broda iznad Kitsa.

Konzul je svirao klavir.

Nekada je taj instrument bio deo luksuza diplomate Hegemonije.

Sada je bio samo klavir na brodu jednog čoveka u svetu bez Hegemonije.

Svirao je dugo.

Ostali su pili vino.

Gledali svetla grada.

Razgovarali o ljudima koji više nisu bili sa njima.

\bigskip

Pitali su se za Sola i Rejčel.

Niko nije znao gde su završili.

Pitali su se za Direa.

Poslednje poruke govorile su o njegovoj novoj ulozi u Crkvi.

Setili su se Kasada.

Njegova sudbina više nije bila nepoznata.

Počivao je u Kristalnom Monolitu, vekovima udaljen od bitke u kojoj je poginuo.

Ako Dire ponovo umre, dve kruciforme možda će još jednom promeniti odnos između dvojice sveštenika.

Niko nije znao.

\bigskip

Teo je gledao grad.

``Šta misliš da ćemo pronaći kada se vratimo među stare svetove Mreže?''

Konzul je prestao da svira.

``Haos.''

``Samo to?''

``Glad na nekim mestima. Ratove na drugima. Nove vlade. Stare diktatore. Ljude koji pokušavaju da ponovo naprave ono što su izgubili.''

``Zvuči obećavajuće.''

Konzul se nasmešio.

``Ali većina će preživeti.''

``Kako znaš?''

``Ne znam.''

Martin je podigao čašu.

``Drugi dobar naučni princip večeras.''

\bigskip

Bron je ćutala.

Ruka joj je počivala na stomaku.

Negde u njenom detetu bila je budućnost o kojoj joj je Kits govorio.

Ona Koja Podučava.

Nije znala šta taj naziv znači.

Nije bila sigurna ni da želi da zna.

Za sada joj je bilo dovoljno da dete bude zdravo.

\bigskip

Sledećeg jutra Bron je trebalo da ode.

Put do Lususa bez dalekobacača više nije bio jednostavan.

Putovanja su ponovo trajala.

Vremenski dug se vratio kao svakodnevna činjenica ljudskog života.

Ali želela je da se vrati kući.

I da tamo rodi Džonijevu kćer.

\bigskip

Pre polaska još jednom je srela Kitsa.

Rekao joj je da Srž više nema prostor u kojem je nekada živela.

Njegova ličnost moraće da pronađe drugi dom.

Bron mu je ponudila ideju.

On se nasmešio.

Zatim je nestao.

\bigskip

Kasnije, u svemirskoj luci, Konzul je gunđao na brodski računar.

``Nešto nije u redu sa njim.''

``Šta?'' pitala je Bron.

``Tražim proveru sistema, a on mi odgovara stihovima.''

Martin je podigao obrve.

``Kakvim stihovima?''

Konzul je pustio snimak.

Bron se nasmešila.

Nije objasnila zašto.

\bigskip

Možda Džon Kits ipak nije otišao tako daleko.

\bigskip

Brod se pripremao za polazak.

Iza njih je ostajao Hiperion.

Vremenske grobnice više nisu bile zatvorene.

Stara Hegemonija nije postojala.

TehnoCentar je nestao iz svog nekadašnjeg doma.

Proterani više nisu bili samo neprijatelji sa granice.

Ljudski svetovi ponovo su bili udaljeni jedni od drugih.

Vreme i prostor vratili su cenu koju im je Mreža vekovima oduzimala.

\bigskip

Ali čovečanstvo nije nestalo.

Po prvi put posle mnogo vekova nije postojala jedna Mreža koja će svima određivati isti ritam.

Svetovi su ponovo morali da pronađu sopstvene puteve.

Neki će propasti.

Neki će se promeniti.

Neki će izgraditi nešto novo.

\bigskip

A na Hiperionu su i dalje stajale Vremenske grobnice.

Otvorene prema budućnosti.

Čekajući ljude koji još nisu bili rođeni.

I priče koje još nisu počele.